\documentclass[a4paper]{article}
\usepackage[pdftex]{graphicx}
\usepackage[utf8]{inputenc}
\usepackage{enumerate}
\usepackage{siunitx}
\sisetup{locale=DE}
\usepackage{href-ul}
\hypersetup{
	colorlinks=true,
	linkcolor=blue,
	urlcolor=blue} 
\usepackage{icomma}
\usepackage{amssymb}
\usepackage{geometry}
\geometry{a4paper, top=15mm, left=15mm, right=15mm, bottom=15mm,
	headsep=10mm, footskip=12mm}

\begin{document}
{\bf Beschleunigte Bewegung}
\begin{enumerate}[1.]

%Aufgabe 1
{\bf \item Bei der Fußballweltmeisterschaft erreichte der Ball bei einem Elfmeter eine Geschwindigkeit von 108 km/h.} Hätte der Torwart den Ball gehalten, dann wäre der Ball innerhalb von 0,04 s auf die Geschwindigkeit Null abgebremst worden. 

\begin{enumerate}[a)]
\item Wie groß wäre dann die mittlere, als konstant angenommene Verzögerung (= negative Beschleunigung) des Balles beim Abbremsen ?
\item Berechne die Kraft, die der Torwart dann für einen Ball der Masse 420 g dafür hätte aufbringen müssen.
\end{enumerate}

%Aufgabe 2
{\bf \item Aufgabe}
\begin{enumerate}[a)]
\begin{minipage}{0.4\textwidth}
\item Erkläre in Stichworten den Bewegungsverlauf des nebenstehenden t--v--Diagramms.
\item Begründe aus dem Diagramm, dass der in der Phase 1 zurückgelegte Weg doppelt so groß ist wie der in Phase 2.
\end{minipage}
\hfill
\begin{minipage}{0.45\textwidth}
\includegraphics[width=5.5 cm]{vaute104.pdf}
\end{minipage}
\item Berechne den Gesamtweg, dem der Körper in 12 s zurücklegt 
(Falsches Ersatzergebnis: 60 m)\\
Welche Durchschnittsgeschwindigkeit hat er?
\item Der Körper, dessen Bewegung das Diagramm beschreibt, hat eine Masse von 750 g. In welcher der 5 Bewegungsphasen ist der Betrag der einwir\-kenden Kraft am größten bzw. am kleinsten? Wie groß sind die Kräfte jeweils?
\item Der Körper fährt nun sofort (bei t = 12 s) den gesamten Weg zurück. er startet dazu mit der Beschleunigung $|\vec{a}|=\SI{4,0}{\meter\over {\second^2}}$ und legt auf diese Weise 32 m zurück. Die restliche Strecke bewältigt er dann mit unveränderter Geschwindigkeit. Wie lange braucht er für die Gesamtstrecke zurück?
\item Zeichne qualitativ das t--v--Diagramm dieser Rückfahrt.
\end{enumerate}

{\bf \item Ein Auto fährt mit konstanter Geschwindigkeit auf der Autobahn, der Tacho zeigt 
 $\SI{125}{\kilo\meter\over \hour}$.} Bei einer Messung wird für eine Strecke von $s_1=\SI{5,00}{\kilo\meter}$ die Zeit 
$t= \SI{2}{\minute}~\SI{30}{\second}$ gestoppt. 
\begin{enumerate}[a)]
\item Berechne die tatsächlich gefahrene Geschwindigkeit $v$ in $ \SI{}{\meter\over \second} $ und 
 $ \SI{}{\kilo\meter\over \hour} $
\item Welche Zeit benötigt das Auto bei dieser Geschwindigkeit für eine $s_2=\SI{5,40}{\kilo\meter}$
lange Strecke? Verwende bei der Antwort die Einheiten $\SI{}{\minute}$ und $\SI{}{\second}$.
\item Welche Strecke hätte das Auto in $\SI{1,6}{\hour}$ zurückgelegt, wenn die Tachoanzeige richtig wäre? 
\item Formuliere eine gleichwertige Aussage dafür, dass es sich um eine Bewegung mit konstanter Geschwindgkeit handelt. 

\end{enumerate}

\begin{minipage}{0.65\textwidth}
%Aufgabe 1
{\bf \item Ein Vergleich eines Motorsportmagazins ergab, dass von 7 getesteten Supersportwagen der Porsche 911 Turbo S die besten Beschleunigungswerte besitzt.} Dieses Auto beschleunigt von 0 auf $ \SI{200}{\kilo\meter\over \hour}$ in genau 
$\SI{10,0}{\second}$.
\end{minipage}
\hfill
\begin{minipage}{0.3\textwidth}
\includegraphics[width=4 cm]{porsche129}
\end{minipage}

\begin{enumerate}[a)]
\item Berechne die Beschleunigung $a$ des Porsches. Dabei soll angenommen werden, dass die Beschleunigung $ a $ konstant ist. 
\item Ermittle die Strecke, die der Wagen innerhalb dieser $\SI{10}{\second}$ zurücklegt.
\end{enumerate}


{\bf Bei einer weiteren Testfahrt wird der Porsche aus dem Stillstand $(t_0=\SI{0}{\second})$ auf $\SI{60}{\meter\over \second} $ innerhalb von $ (t_1=\SI{15}{\second})$ beschleunigt.} Ab  $ t_1=\SI{15}{\second}$ behält der Sportwagen konstant diese Geschwindigkeit bis  $ t_2=\SI{45}{\second}$ bei, bevor er danach innerhalb von 20 Sekunden mit konstanter Beschleunigung auf $\SI{90}{\meter\over \second} \quad   (t_3=\SI{65}{\second})$ beschleunigt.

\begin{enumerate}[(i)]
\item Zeichne ein t-v-Diagramm vom Zeitpunkt $t_0$ bis zum Zeitpunkt $t_3$ . Verwende dabei einen geeigneten Maßstab.
\item Bestimme mit Hilfe des Diagramms den insgesamt zurückgelegten Weg.
\end{enumerate}

\end{enumerate} 

\fbox{
	\begin{minipage}{0.5\textwidth}
		Zur Lösung bitte \href{https://www.okuyakl.de/phys/p10kinL410/ll410.pdf}{hier klicken} oder den QR-Code scannen.\\
	Weitere Arbeitsblätter gibt es unter 
	
	\href{https://www.okuyakl.de}{www.okuyakl.de}
	\end{minipage}
	\hfill
	\begin{minipage}{0.4\textwidth}
		\includegraphics[width=1.5 cm]{../../viecher/zwe03}
		\includegraphics[width=3 cm]{qr410}
		\includegraphics[width=2 cm]{../../viecher/afanticon1}
		
	\end{minipage}}

\end{document}
